% Скриншоты UI: каталог docs/diplom/фото/ (подпапки по экранам).
% Список проектов: ui.projects.list/1.jpg  |  Карточка проекта: ui.projects/1.jpg, 2.jpg
\subsection{Экраны веб-интерфейса \texttt{Chronos.Master.Web}}

Ниже приведены скриншоты экранов в соответствии с маршрутами React (\texttt{App.tsx}). На рисунки даны отсылки в тексте главы~3 и настоящей главы.

\begin{figure}[!ht]
  \centering
  \IfFileExists{фото/ui.dashboard/1.png}{%
    \includegraphics[width=0.94\linewidth]{ui.dashboard/1.png}%
  }{%
    \IfFileExists{фото/ui.dashboard/1.jpg}{%
      \includegraphics[width=0.94\linewidth]{ui.dashboard/1.jpg}%
    }{%
      \begin{minipage}{0.94\linewidth}\centering
        \fcolorbox{black}{gray!8}{%
          \begin{minipage}{0.92\linewidth}\centering
            \vspace{1.2cm}
            {\small\textit{Скриншот дашборда (\texttt{/ui/dashboard}): положите кадр в \texttt{фото/ui.dashboard/1.png} или \texttt{1.jpg}.}}
            \vspace{1.2cm}
          \end{minipage}}
      \end{minipage}%
    }%
  }
  \caption{Главная панель (дашборд) веб-интерфейса}
  \label{fig:ui-dashboard}
\end{figure}

\begin{figure}[!ht]
  \centering
  \IfFileExists{фото/ui.agents/1.jpg}{%
    \includegraphics[width=0.94\linewidth]{ui.agents/1.jpg}%
  }{%
    \begin{minipage}{0.94\linewidth}\centering
      \fcolorbox{black}{gray!8}{%
        \begin{minipage}{0.92\linewidth}\centering
          \vspace{1.2cm}
          {\small\textit{Скриншот списка агентов (\texttt{/ui/agents}): положите кадр в \texttt{фото/ui.agents/1.jpg}.}}
          \vspace{1.2cm}
        \end{minipage}}
    \end{minipage}%
  }
  \caption{Список агентных узлов}
  \label{fig:ui-agents}
\end{figure}

\begin{figure}[!ht]
  \centering
  \IfFileExists{фото/ui.projects.list/1.jpg}{%
    \includegraphics[width=0.94\linewidth]{ui.projects.list/1.jpg}%
  }{%
    \begin{minipage}{0.94\linewidth}\centering
      \fcolorbox{black}{gray!8}{%
        \begin{minipage}{0.92\linewidth}\centering
          \vspace{1.2cm}
          {\small\textit{Положите скриншот списка проектов в файл \texttt{фото/ui.projects.list/1.jpg} (отдельно от кадров карточки проекта в \texttt{ui.projects/}).}}
          \vspace{1.2cm}
        \end{minipage}}
    \end{minipage}%
  }
  \caption{Список compose-проектов кластера}
  \label{fig:ui-projects}
\end{figure}

\begin{figure}[!ht]
  \centering
  \includegraphics[width=0.94\linewidth]{ui.projects/1.jpg}\\[2mm]
  \includegraphics[width=0.94\linewidth]{ui.projects/2.jpg}
  \caption{Детальная страница проекта}
  \label{fig:ui-project-detail}
\end{figure}

\begin{figure}[!ht]
  \centering
  \includegraphics[width=0.94\linewidth]{ui.projects.archive/1.jpg}
  \caption{Архивированные проекты}
  \label{fig:ui-archived}
\end{figure}

\begin{figure}[!ht]
  \centering
  \includegraphics[width=0.94\linewidth]{ui.routing/1.jpg}\\[2mm]
  \includegraphics[width=0.94\linewidth]{ui.routing/2.jpg}
  \caption{Управление TCP-маршрутизацией (HAProxy)}
  \label{fig:ui-routing}
\end{figure}

\begin{figure}[!ht]
  \centering
  \includegraphics[width=0.94\linewidth]{ui.sandbox/1.jpg}\\[2mm]
  \includegraphics[width=0.94\linewidth]{ui.sandbox/2.jpg}\\[2mm]
  \includegraphics[width=0.94\linewidth]{ui.sandbox/3.jpg}
  \caption{Песочница YAML и Fluent API}
  \label{fig:ui-sandbox}
\end{figure}
